% Design notes: site-first enumeration with asymmetric displacements.
%
% A working document, not a paper.  Records the decisions taken, the scheme
% they define, the parameters still to be chosen, and the order the work
% should be done in.  Companion to gbShiftsEnumeration.tex (the scheme this
% replaces), plasticDisplacement.tex and mesostateOverview.tex.

\documentclass[11pt]{article}
\usepackage{amsmath}
\usepackage{amssymb}
\usepackage[margin=1.1in]{geometry}
\usepackage{enumitem}
\setlist[itemize]{itemsep=2pt, topsep=4pt, parsep=0pt}
\setlist[enumerate]{itemsep=3pt, topsep=4pt, parsep=0pt}

\newcommand{\A}{\mathcal{A}}
\newcommand{\B}{\mathcal{B}}
\newcommand{\D}{\mathcal{D}}
\newcommand{\nhat}{\hat{\mathbf{n}}}
\newcommand{\dmax}{d_{\max}}
\newcommand{\dmin}{\delta_{\min}}

\title{Site-first enumeration with asymmetric displacements\\
       \large design notes}
\date{}

\begin{document}
\maketitle

\section*{Decisions taken}

\begin{enumerate}
  \item \textbf{Enumerate coincidence sites first}, then the atoms that can meet
        at each, rather than enumerating translations and deriving the sites.
  \item \textbf{Allow the two grains to be displaced independently.} The
        symmetric split $\pm\mathbf{t}/2$ becomes one case among many rather
        than a constraint.
  \item \textbf{Sites are taken from $\tfrac{1}{2}\D$}, the lattice of midpoints
        of $\A$ and $\B$ atoms, with an optional filter restricting them to
        $\D$.
\end{enumerate}

The third decision was reversed once during the discussion and is worth
recording. $\D$ sites have the property that both nodal displacements are
themselves DSC vectors, which $\tfrac{1}{2}\D$ sites do not. But the midpoint of
an $\A$ atom and a $\B$ atom lies on $\D$ only when the translation is an
\emph{even} DSC vector --- an index-8 sublattice --- so $\D$ sites reproduce
only about one flat state in eight exactly. Losing the flat baseline costs more
than the crystallographic tidiness gains: the jump
$\mathbf{t} = \mathbf{t}_\A - \mathbf{t}_\B$ is a DSC vector either way, and it
is the jump that carries the physics. Since $\D \subset \tfrac{1}{2}\D$,
enumerating the larger set and filtering forecloses nothing.

\section*{The scheme}

\subsection*{The node record}

A node is a triple $(\mathbf{s}, \mathbf{t}_\A, \mathbf{t}_\B)$, read as the two
records $(\mathbf{t}_\A, \mathbf{s})$ and $(\mathbf{t}_\B, \mathbf{s})$ sharing a
coincidence point:
\begin{align*}
  \mathbf{x}_\A &= \mathbf{s} - \mathbf{t}_\A \in \A,
  &\mathbf{u}_\A &= \mathbf{t}_\A, \\
  \mathbf{x}_\B &= \mathbf{s} - \mathbf{t}_\B \in \B,
  &\mathbf{u}_\B &= \mathbf{t}_\B,
\end{align*}
with the jump $\mathbf{t} = \mathbf{t}_\A - \mathbf{t}_\B
= \mathbf{x}_\B - \mathbf{x}_\A \in \D$ derived rather than stored. The present
scheme is the special case $\mathbf{t}_\A = +\mathbf{t}/2$,
$\mathbf{t}_\B = -\mathbf{t}/2$.

\subsection*{Enumeration}

\begin{enumerate}
  \item Sites: all points of $\tfrac{1}{2}\D$ in a slab of half-thickness $T$
        about the flat boundary, taken over one in-plane period cell.
  \item For each site, the $\A$ atoms with
        $\lVert \mathbf{s}-\mathbf{x}_\A \rVert \le \dmax$ and the $\B$ atoms
        with $\lVert \mathbf{s}-\mathbf{x}_\B \rVert \le \dmax$.
  \item Every combination of one from each is a candidate node. A site with
        $n_\A$ and $n_\B$ atoms in range yields $n_\A n_\B$ nodes.
\end{enumerate}
The bound is on how far \emph{each atom} moves, independently; it implies
$\lVert\mathbf{t}\rVert \le 2\dmax$ but is not equivalent to it.

\subsection*{Validity of a mesostate}

With $\Lambda$ the bicrystal box lattice, $L = \langle \mathbf{p}_1,
\mathbf{p}_2 \rangle$ the in-plane sublattice and $\Pi$ the projection onto the
boundary plane, a set of nodes is admissible when, pairwise:
\begin{itemize}
  \item $\mathbf{x}_\A \bmod \Lambda$ are distinct --- no two nodes on one $\A$ atom;
  \item $\mathbf{x}_\B \bmod \Lambda$ are distinct --- no two nodes on one $\B$ atom;
  \item $d_L(\mathbf{s},\mathbf{s}') \ge \dmin$, where $d_L$ is the distance
        between the projected sites on the torus $\text{plane}/L$.
\end{itemize}
The third subsumes exact coincidence ($d_L = 0$) and so replaces, rather than
supplements, a separate projected-site key. It is what removes the
mesh-collapse failures that reject $2.5\%$ of states under the present scheme.

All three are pairwise and closed under taking subsets, so admissibility is
hereditary and can be enforced \emph{during} enumeration: a depth-first walk
emits a state at every node of its recursion tree and discards nothing.

\subsection*{Mechanics}

Precompute per node its two atom classes and its site class; precompute per site
the closed neighbourhood $N[\sigma] = \{\sigma\} \cup \{\sigma' : d_L < \dmin\}$.
Gauss-reduce $(\mathbf{p}_1,\mathbf{p}_2)$ before any nearest-lattice-point
rounding, or an oblique cell gives wrong minima. The walk then carries two
boolean arrays over atom classes and an integer \emph{counts} array over sites,
incremented over $N[\sigma]$ on descent and decremented on return.

Ceiling on the size of a state:
\[
  \min\bigl(\,\lvert \A/\Lambda \rvert,\;
             \lvert \B/\Lambda \rvert,\;
             \text{packing number of the torus at radius } \dmin/2 \,\bigr).
\]

\section*{Parameters to be chosen}

Three parameters replace \texttt{tMax}, \texttt{tPerpMax} and
\texttt{sPerpMax}. \texttt{tPerpMax} disappears: it restricted the direction of
the translation, which has no role once the region bounds the site.

\begin{itemize}
  \item $T$ --- \textbf{slab half-thickness}, the region about the flat boundary
        from which sites are drawn. Controls how far the boundary may facet out
        of plane. Site count grows linearly in $T$.
  \item $\dmax$ --- \textbf{how far one atom may move} to reach its site.
        Controls how many atoms are in range of each site, so cost grows as
        $\dmax^3$ per side and $\dmax^6$ in the node count. The dominant cost
        driver.
  \item $\dmin$ --- \textbf{minimum in-plane separation of coincidence points.}
        Physically bounded below by the scale at which two coincidence points
        cease to be distinguishable structures --- a fraction of the
        nearest-neighbour spacing --- and above by the in-plane site spacing,
        past which legitimate distinct sites start being excluded wholesale.
\end{itemize}
\texttt{maxEngaged} and \texttt{maxStates} survive unchanged but stop being
safety rails: they become the primary controls on the size of a sweep.

\subsection*{Estimating the cost before choosing}

For the $\Sigma 5$ case at hand, with $a_0 = 3.615$\,\AA{} (so
$\rho_\A = \rho_\B = 0.0847$\,\AA$^{-3}$ and $b = 2.556$\,\AA), an in-plane cell
area of $29.2$\,\AA$^2$, and $\lvert \D : \A \rvert = \Sigma$:
\[
  \rho_{\D} = \Sigma\,\rho_\A \approx 0.42\ \text{\AA}^{-3},
  \qquad
  \rho_{\frac{1}{2}\D} = 8\,\rho_{\D} \approx 3.4\ \text{\AA}^{-3} .
\]
Hence, per in-plane cell,
\[
  N_{\text{sites}} \approx \rho_{\frac{1}{2}\D}\times 29.2 \times 2T
  \approx 99\,(2T)\ \text{sites per \AA{} of slab},
  \qquad
  n_\A(\dmax) \approx \tfrac{4}{3}\pi\,\rho_\A\,\dmax^{3}
  \approx 0.35\,\dmax^{3},
\]
and the node count is $N_{\text{sites}}\,n_\A n_\B$. Some indicative
combinations:
\begin{itemize}
  \item $2T = 0.5$\,\AA, $\dmax = 1.5$\,\AA: $\approx 50$ sites, $\approx 1.2$
        atoms in range per side, order $70$ nodes.
  \item $2T = 1$\,\AA, $\dmax = 2.0$\,\AA: $\approx 99$ sites, $\approx 2.8$ per
        side, order $800$ nodes.
  \item $2T = 2$\,\AA, $\dmax = 2.56$\,\AA: $\approx 198$ sites, $\approx 5.9$
        per side, order $7000$ nodes.
\end{itemize}
These are averages over a uniform density and will be wrong in detail --- the
site and atom lattices are structured, not random --- so they indicate the
regime, not the answer.

\subsection*{Reproducing the flat states}

To contain every state the present enumeration can build, take
\[
  \dmax \ge \tfrac{1}{2}\,\texttt{tMax}\cdot b,
  \qquad
  T \ge \tfrac{1}{2}\,\texttt{sPerpMax}\cdot b .
\]
The first because the symmetric split moves each atom by half the translation.
Note that even so the containment is at the \emph{node} level only: any
$\dmin > 0$ removes states whose engaged sites are closer than the threshold,
which is deliberate but means the new state set is not a superset of the old.

\section*{Order of work}

The parameters should be chosen from measured counts, not from the estimates
above, which argues for building the enumerator in two stages.

\begin{enumerate}
  \item \textbf{Counting probe.} Enumerate sites and nodes and report, for
        candidate $(T, \dmax, \dmin)$: the site count, the node count, the three
        class counts, and the resulting ceiling. No mesostates constructed, no
        file format committed. This is what the parameters should be chosen
        from.
  \item \textbf{The record and its format.} Decide the on-disk representation of
        a node. This constrains the enumerator, the driver and the meaning of a
        signature simultaneously, so it should be settled before any of them is
        written.
  \item \textbf{The enumerator proper}, replacing the translation-first search.
  \item \textbf{\texttt{getFacetedSurfaces()}}, which is the only place the
        symmetric split is hardcoded; it takes the triple directly.
  \item \textbf{The driver}, whose site-grouping and clash handling become the
        three-key walk described above.
\end{enumerate}

Untouched throughout: \texttt{GbContinuum}, \texttt{GbFacet},
\texttt{assertFacetsGlue()} and \texttt{gluedCloudB}. All of these already take
the two node lists independently and nowhere assume the displacements are
opposite, so asymmetric splits need no change below the mesostate level.

\section*{Consequences to expect}

\begin{itemize}
  \item Node counts of order $10^2$ to $10^4$ against the present $42$, and
        exhaustive sweeps correspondingly out of reach: \texttt{maxEngaged}
        carries the load.
  \item The $62{,}399$-state sweep ceases to be a state-by-state baseline. Its
        energies were computed with the superseded constant-per-face quadrature
        in any case, so they need recomputing regardless; what matters is that
        the flat reference can be rebuilt on the same footing as the faceted
        states it is compared against, which the $\tfrac{1}{2}\D$ choice
        preserves.
  \item The mesh-collapse rejections disappear from the construction stage,
        being excluded during enumeration instead.
\end{itemize}

\section*{Still open}

\begin{enumerate}
  \item Whether the slab bounds the site $\mathbf{s}$ or the reference atom
        $\mathbf{x}_\A$. They differ by at most $\dmax$. Bounding the site
        controls the deformed boundary, which is what ``faceted'' refers to;
        bounding the atom is the more literal reading of ``atoms near the
        boundary''.
  \item Whether $\dmax$ bounds each displacement separately (assumed here) or
        their sum.
  \item Numerical values of $T$, $\dmax$, $\dmin$ --- to come from the counting
        probe.
  \item The on-disk record format.
\end{enumerate}

\end{document}
